\chapter{AI-Based Safety-Critical Systems}
Although inserting AI into a safety-critical system can be very challenging, there are many techniques that can be used and have already been tested. For example, referring to the classification in Section~\ref{sec:application}, we find that \textit{Class I} and \textit{Class II} include systems that are already on the market (e.g., SIL4 railway interlocking).

%%%%%%%%%%%%%%%%%%%%%%%%%%%%%%%%%%%%%%%%%%%%%%%%
\section{Product}
As mentioned in Section~\ref{sec:application}, the product class describes a safety-critical system that relies on its SIS on one or more AI-based functions trained offline. These functions can be "simple", as in the case of a symbolist AI model, thus belonging to \textit{Class I}, or they can be very complex and belong to \textit{Class II}. In the first case, since the model is deterministic, we can use formal verification to ensure the system's compliance with FuSa. In the second case, the standard industry technique is to use a runtime checker that makes a safe choice for the next action if the model operates outside its requirements or produces an unsafe result.

%%%%%%%%%%%%%
\subsection{Testing}
The main problem with these AI-based safety-critical systems is the assurance of safety that we can assign to them. This stems from the stochastic nature of AI models, which implies that they can have \textit{uncertainty failures}. To reduce this uncertainty, we must investigate unknown factors. Therefore, when describing safety functions in this scenario, we speak about intended functionality: their designs are a set of high-level objectives rather than a set of precise and well-defined instructions. This can be mitigated by applying formal verification methods, but even in this case, the black-box nature of the models suggests that these methods cannot be the only ones. Another challenge is the testing of the AI model itself, which is a widely explored area (e.g., test generation from an NLP model).

%%%%%%%%%%%%%
\subsection{Neural Networks}
Although symbolist models provide understandability and explainability for their decisions, the most powerful models are Neural Networks (NNs).

When dealing with this kind of model, we must consider several detailed issues:
\begin{itemize}
    \item Supervised learning relies on data. When managing data (i.e., preprocessing and dividing data into training and test sets), we must do so with very high accuracy.
    \item This is an area where \textit{changing nothing changes everything}~\cite{perez2024artificial}. This means that changing activation functions and hyperparameters can lead to very different behaviors.
    \item NN training is iterative; this must be included in the V-Model within the specification and testing phases.
    \item AI frameworks (e.g., PyTorch) are general-purpose and user-friendly. However, there are some APIs that also provide a set of safety requirements.
\end{itemize}

%%%%%%%%%%%%%%%%%%%%%%%%%%%%%%%%%%%%%%%%%%%%%%%%
\section{Runtime}
As discussed in Section~\ref{sec:application}, a runtime safety-critical system is able to change its AI module during execution. These are notable systems because they allow the SIS to evolve with the context: a case study reported a gas oil turbine where the safety function evolves as the engine degrades.

This kind of solution bases its dependability and standard compliance on four different techniques:
\begin{itemize}
    \item \textbf{Safety Bag} \\
        This is the same technique mentioned above: a runtime checker monitors the AI model's results, and if it detects strange behavior, it provides a safe result.
    \item \textbf{Safe Adaptation} \\
        Both the AI model and the (continuous) training algorithm conform to safety standards.
    \item \textbf{Limited Adaptation} \\
        The (continuous) training is limited within a safe range: this means that even if the weights of the model change, they cannot do so with high magnitude.
    \item \textbf{Limited Actuation} \\
        In this case, the adaptation is performed with hard constraints on the resulting model's output: for example, the new model cannot have an output where the magnitude exceeds a certain threshold.
    \item \textbf{Library-Based Offline} \\
        We train a set of AI models, each of which conforms to safety standards. During inference, the system chooses from all possible models in a non-linear way (i.e., with \texttt{if/else}).
\end{itemize}

%%%%%%%%%%%%%%%%%%%%%%%%%%%%%%%%%%%%%%%%%%%%%%%%
\section{Process}
These are offline solutions that help and support engineers in crafting safety-critical systems, both from the perspective of the overall system and the AI component that implements the safety function.

\myskip

In traditional safety engineering, AI models can be useful for validating and assessing the V-Model. During specification, we can use an NLP model (e.g., transformer-based) to discover hazards and issues in the process; during design and implementation, models can be useful for generating code; and during the testing phase, they become useful for generating test cases.

If we look at the use of AI to construct other models that will implement safety functions, there are different solutions: in AutoML, we use reinforcement learning to find the best model hyperparameters, and we can implement models that search in an exhaustive way for all combinations to test our system.